% Source: https://artofproblemsolving.com/community/c3h1566509_mathcounts_mock_test_2015
% Stable identity: aops-topic-1566509-q26
Let $M$ be the midpoint of $AC$. The altitude $BM$ is perpendicular to $AC$, with $AM=MC=9$ and $BM=\sqrt{15^2-9^2}=12$. As the triangle turns about $AC$, the point $B$ traces a circle of radius $BM=12$ centred at $M$ in the plane perpendicular to $AC$. The solid is two cones glued along this circular base: each has radius $12$ and height $9$ (from $M$ to $A$, and from $M$ to $C$).

\begin{asy}
usepackage("amsmath");
size(6cm);
defaultpen(fontsize(10pt));
pair A=(-9,0),C=(9,0),B=(0,12),Bp=(0,-12);
draw(A--B--C,linewidth(1));draw(A--Bp--C,gray+dashed);
draw(A--C,dashed+red);
path el=shift(0,0)*xscale(2.4)*yscale(12)*unitcircle;
draw(subpath(el,1,3),gray+dashed);draw(subpath(el,3,5),gray);
draw(B--(0,0),blue);
real q=0.8;draw((0,q)--(q,q)--(q,0),blue);
label("$A$",A,W);label("$C$",C,E);label("$B$",B,N);label("$B'$",Bp,S);
label("$15$",(A+B)/2,NW);label("$15$",(B+C)/2,NE);
label("$9$",(-4.5,0),S);label("$9$",(4.5,0),S);label("$12$",(0,4),E,blue);
label("\mbox{}",(min(currentpicture,user=true)+max(currentpicture,user=true))/2,black); // keeps the text after the figure black when LaTeX typesets the labels inline
\end{asy}

(The axis $AC$ is red; $B'$ is the position of $B$ after half a turn, and the ellipse is the common base seen at an angle.) Hence
\[
V=2\cdot\dfrac{1}{3}\pi\cdot12^2\cdot9=\boxed{864\pi}.
\]
So the radius is $12$ but the height of each cone is $9$, not $12$: the axis of rotation is $AC$, so the heights are measured along $AC$.
